\section{Requirements and Formal Specification}
\label{sec:spec}

\subsection{Users and requirements}

Three user roles drive the requirements. \emph{Model researchers} add and
compare estimators without rebuilding dataset parsing, splits, metrics, or
provenance capture. \emph{Dataset maintainers} publish adapters for their
archives with source identity, licence, checksums, and recording boundaries
intact. \emph{Reviewers and reproducibility researchers} must be able to
verify---without trusting the authors---that recordings do not cross
partitions, that held-out conditions did not select models, that every
registered seed is accounted for, and that results map to specific data,
code, and environment artefacts. All three roles interact with the framework
only through its public contracts; the requirements therefore reduce to one
sentence: \emph{extensions must be expressive enough for heterogeneous
datasets and models, yet structurally incapable of violating experiment
integrity}.

\subsection{Experiment lifecycle}

We model one (task, model-plugin) experiment as a guarded state machine over
the states
\textsc{Empty} $\rightarrow$ \textsc{SourcesVerified} $\rightarrow$
\textsc{TaskAudited} $\rightarrow$ \textsc{DataPrepared} $\rightarrow$
\textsc{ModelConstructed} $\rightarrow$ \textsc{Trained}$_s$ $\rightarrow$
\textsc{Evaluated}$_s$ $\rightarrow$ \textsc{ResultValidated} $\rightarrow$
\textsc{Published}, plus a terminal \textsc{Rejected} state. Each transition
$t_1$\,--\,$t_8$ carries a guard; a guard that fails sends the experiment to
\textsc{Rejected} with a machine-readable reason, and nothing partial
survives: the discipline is fail-closed. Table~\ref{tab:transitions} lists
the transitions with their guards and the runtime mechanism that enforces
each.

\begin{table}[t]
\centering\small
\caption{Lifecycle transitions, guards, and enforcing mechanisms.}
\label{tab:transitions}
\begin{tabular}{@{}llp{0.42\linewidth}@{}}
\toprule
Transition & Guard (invariant) & Runtime mechanism \\
\midrule
$t_1$ acquire & INV-1 & checksum-verified download; re-verification at every stream open \\
$t_2$ audit & INV-2 (recordings) & fail-closed recording audit \\
$t_3$ prepare & INV-2 (windows), INV-3 & window-inheritance check; kernel-restricted fit scope \\
$t_4$ construct & INV-6 & plugin registry and contract validation \\
$t_5$ train & INV-3, INV-4 & partition-restricted data views; versioned training policy \\
$t_6$ evaluate & INV-3, INV-4 & test view materialized only here, once per seed \\
$t_7$ validate & INV-4, INV-5 & seed-bijection and provenance validation; input-hash recompute \\
$t_8$ publish & INV-1, INV-5 & immutable release with checksums \\
\bottomrule
\end{tabular}
\end{table}

\subsection{Integrity invariants}

Six invariants hold in every committed state:

\begin{description}
\item[INV-1 (Source integrity).] Every data file matches its registered byte
size and SHA-256, and every result references its input artefacts by content
hash.
\item[INV-2 (Partition exclusivity).] The partition assignment is a total
function over recordings; every window inherits its recording's partition;
no recording or window appears in two partitions.
\item[INV-3 (Test isolation).] Test data is read only in $t_6$; preprocessing
statistics are fitted on train/validation only; model selection and early
stopping use validation only. Enforced structurally: training plugins never
receive test rows.
\item[INV-4 (Seed completeness).] Result runs biject with the task's
registered seeds, in the frozen order; failures are retained as explicit
placeholders.
\item[INV-5 (Provenance completeness).] A valid result names its task,
dataset, plugin, code commit, environment, and training policy, and its
input hashes recompute.
\item[INV-6 (Plugin contract).] External plugins satisfy interface, identity,
and version constraints; collisions and mismatches fail closed.
\end{description}

\subsection{Model checking}

The lifecycle and invariants are written in TLA\textsuperscript{+}
(\texttt{formal/ExperimentLifecycle.tla}). Bad candidates---tampered
partitions, windows claiming a foreign partition, fits that attempt test
data, contract-less plugins, incomplete bundles---are explored
non-deterministically, so every guard's rejection path is exercised, not
assumed. TLC checks nine properties (type safety, the six invariants, and
the fail-closed discipline) over a small finite model (three recordings,
three windows, two seeds). The result: \emph{model checking completed with
no errors}---89{,}476 states generated, 46{,}686 distinct states, search
depth 11, no deadlocks. Two modelling defects were found and fixed by TLC
during development (an invariant scoped to a state where its variables are
not yet meaningful, and an underspecified rejecting branch), which is
evidence the check has teeth. The archived log and a reproduction script
ship with the repository, and CI re-runs the check on every push.

\paragraph{Scope of the formal claim.}
This is bounded model checking of an \emph{abstract} lifecycle: it verifies
the guard logic of the design. It is not a proof about the Python
implementation. Implementation-level coverage comes from the runtime
contract checks (\S\ref{sec:impl}), the property-based tests, and the
fault-injection catalogue (\S\ref{sec:eval}); we keep these layers explicit
rather than letting the phrase ``formally specified'' absorb them.
